\documentclass{article}
\usepackage[utf8]{inputenc}
\usepackage[T1]{fontenc}
\usepackage{geometry}
\geometry{
    a4paper,
    left=2.5cm,
    right=2.5cm,
    top=2.5cm,
    bottom=2.5cm
}
\usepackage{graphicx}
\usepackage{float}
\usepackage{hyperref}
\usepackage[english]{babel}
\usepackage{amsmath}
\usepackage{amssymb}
\usepackage{array}
\usepackage{booktabs}
\usepackage{caption}
\usepackage{subcaption}
\usepackage{enumitem}
\usepackage{fancyhdr}
\usepackage{titlesec}
\usepackage{url}
\usepackage{textcomp}
\usepackage{xcolor}
\usepackage{listings}
\usepackage{multirow}
\usepackage{longtable}
\usepackage{tabularx}
\usepackage{times}

% Set up page style
\pagestyle{fancy}
\fancyhf{}
\rfoot{\thepage}
\renewcommand{\headrulewidth}{0pt}

% Title formatting
\titleformat{\chapter}[display]{\normalfont\bfseries}{\chaptertitle\thechapter}{0pt}{\Huge}
\titlespacing*{\chapter}{0pt}{50pt}{40pt}

% Document info
\title{Template Document}
\author{Academic Support Agent}
\date{\today}

\begin{document}
\maketitle

% Seleciona o idioma do documento
	\selectlanguage{brazil}
	
	% Retira espa o extra obsoleto entre as frases.
	\frenchspacing 
	
	% ----------------------------------------------------------
	% ELEMENTOS PR -TEXTUAIS
	% ----------------------------------------------------------
	
	% Se desejar escrever o artigo em duas colunas, descomente a linha abaixo
	% e a linha com o texto "FIM DE ARTIGO EM DUAS COLUNAS".
	% \twocolumn[    		% INICIO DE ARTIGO EM DUAS COLUNAS
    
	% t tulo
	\maketitle
    
	% resumo em portugu s
	\begin{resumoumacoluna}
		
		Neste trabalho   apresentada a formata o que deve ser utilizada nos artigos a serem submetidos ao final dos cursos de Gradua o e P s-Gradua o do IFSP c mpus S o Jo o da Boa Vista. Leia com aten o este documento. O m ximo de palavras para o resumo   250 (duzentos e cinquenta). 
		
		\vspace{\onelineskip}
		
        % todas em letras min sculas, separadas por ponto e v rgula (;)
        \textbf{Palavras-chave}: palavra-chave 1; palavra-chave 2; palavra-chave 3; palavra-chave n.
		
	\end{resumoumacoluna}
	
    % resumo em ingl s
    \renewcommand{\resumoname}{Abstract}
    \begin{resumoumacoluna}
    	
        \begin{otherlanguage*}{english}
   	    	Resumo em ingl s.
   		
       		\vspace{\onelineskip}
   		
   	    	\noindent
   	    	\textbf{Keywords}: keyword 1. keyword 2. keyword 3. keyword n.
   		   
      \end{otherlanguage*} 
  
   	\end{resumoumacoluna}
        
	% ]  				% FIM DE ARTIGO EM DUAS COLUNAS
	
	% ----------------------------------------------------------
	% ELEMENTOS TEXTUAIS
	% ----------------------------------------------------------
	\textual
	
	% ----------------------------------------------------------
	% Introdu o
	% ----------------------------------------------------------
	\section{Introdu o}
	
    Neste documento est o listadas as se es obrigat rias que voc  dever  fornecer, bem como os exemplos dos comandos mais comuns que ser o utilizados na constru o de seu documento. Para pesquisar sobre mais comandos, recomenda-se a utiliza o do site \url{https://ctan.org/}, que   a biblioteca principal do \LaTeX, e o do site \url{https://tex.stackexchange.com} que   uma das principais comunidades para solu o de d vidas relacionadas a \LaTeX. Ambas s o em ingl s.	
    
    
%%%%%%%%%%%%%%%%%%%%%%%%%%%%%%%%%%%%%%%%%%%%%%%%%%%%%%%%%%%%%%%%%%

	\section{Desenvolvimento}
    	
     Texto do desenvolvimento.
     
     
     
     %%%%%%%%%%%%%%%%%%%%%%%%%%%%%%%%%%%%%%%%%%%%%%%%%%%%%%%%%%%%%
     
     \subsection{Revis o da Literatura}
     
     Texto da revis o da literatura.
     
     Este   um exemplo de como usar figuras. Refer ncia cruzada: Figura~\ref{fig:exemplo}
     
     \FloatBarrier
     \begin{figure}[!htbp]
     	\centering
     	\caption{Exemplo de figura}
     	%scale redimensiona a figura.
     	%1.5 = 150% do tamanho original
     	%1 = 100% do tamanho original
     	%0.20 = 20% do tamanho original
     	\includegraphics[scale=0.4]{imagens/exemploFigura}
     	\\\textbf{Fonte:} Elaborada pelo autor
     	\label{fig:exemplo}
     \end{figure}
     \FloatBarrier
     
     
     Este   um exemplo de como usar tabelas. Refer ncia cruzada: Tabela~\ref{tab:exemplo}
     
     \FloatBarrier
     \begin{table}[!htbp]
     	\centering
     	\caption{Exemplo de tabela de 2 colunas}
     	\begin{tabular}{ c | c }
     		\hline
     		\textbf{Coluna 1} & \textbf{Coluna 2} \\ \hline
     		Dado 1a           & Dado 2a           \\ \hline
     		Dado 1b           & Dado 2b           \\ \hline
     		Dado 1c           & Dado 2c           \\ \hline
     		Dado 1d           & Dado 2d           \\ \hline
     	\end{tabular}
     	\\ \vspace{0.2cm}
     	\textbf{Fonte:} Elaborada pelo autor
     	\label{tab:exemplo}
     \end{table}
     \FloatBarrier
     
     
     Este   um exemplo de como usar quadros. Refer ncia cruzada: Quadro~\ref{qua:exemplo}
     
     \FloatBarrier
     \begin{quadro}[!htbp]
     	\centering
     	\caption{Exemplo de quadro}
     	\includegraphics[scale=.7]{imagens/exemploQuadro}
     	\\\textbf{Fonte:} Elaborada pelo autor
     	\label{qua:exemplo}
     \end{quadro}
     \FloatBarrier
     
     
     Este   um exemplo de como usar equa es. Refer ncia cruzada: Equa o~\ref{eq:exemplo}
     
     \begin{equation}
     \sum_{i=1}^{n} i = \frac{n(n+1)}{2}
     \label{eq:exemplo}
     \end{equation}
     
     
     Exemplo de inser o de lista de c digo fonte:
     
     \lstinputlisting[language=Java]{fontes/ClasseExemplo.java} 
     
     
     
     Este   um exemplo de como inserir texto sem formata o (ambiente verbatim):
     
     \begin{verbatim}
     Texto sem formata o, como espa amento igual.
     \end{verbatim}
     
     
     Exemplo de lista de itens:
     
     \begin{itemize}
     	\item \textbf{Item 1:} texto...;
     	\item \textbf{Item 2:} texto...;
     	\begin{itemize}
     		\item \textbf{Subitem:} texto...;
     		\item \textbf{Subitem:} texto...;
     		\item \textbf{Subitem:} texto...;
     	\end{itemize}
     	\item \textbf{Item 3:} texto...;
     	\item \textbf{Item n:} texto....
     \end{itemize}
     
     
     Exemplo de lista numerada:
     
     \begin{enumerate}
     	\item \textbf{Item:} texto...;
     	\item \textbf{Item:} texto...;
     	\begin{enumerate}
     		\item \textbf{Subitem:} texto...;
     		\item \textbf{Subitem:} texto...;
     		\item \textbf{Subitem:} texto...;
     	\end{enumerate}
     	\item \textbf{Item:} texto...;
     	\item \textbf{Item:} texto....
     \end{enumerate}
     
     
     Exemplos de comandos para texto e refer ncias:
     
     \begin{itemize}
     	\item Para iniciar um novo par grafo, basta deixar uma linha em branco no c digo fonte;
     	\item N o force o compilador a pular mais de uma linha, pois ter  influ ncia negativa na composi o do documento;
     	\item Sempre deixe o \LaTeX\ realizar a formata o de par grafos e posicionamento de elementos;
     	\item Utiliza o de aspas simples (abertura \verb|`|, fechamento \verb|'|): `Texto entre aspas simples';
     	\item Utiliza o de aspas duplas (abertura \verb|``|, fechamento \verb|''|): ``Texto entre aspas duplas'';
     	\item Negrito (comando \verb|\textbf|): \textbf{texto em negrito};
     	\item It lico (comando \verb|\textit|): \textit{texto em it lico};
     	\item Sublinhado (comando \verb|\underline|): \underline{texto sublinhado};
     	\item Negrito e it lico (usar comandos juntos): \textbf{\textit{texto em negrito e it lico}};
     	\item Alterar cor do texto (comando \verb|\textcolor{cor}{texto}|):
     	\begin{itemize}
     		\item Exemplo \verb|\textcolor{red}{texto}|: \textcolor{red}{texto vermelho};
     		\item Exemplo \verb|\textcolor[RGB]{255, 102, 0}|: \textcolor[RGB]{255, 102, 0}{texto laranja};
     		\item Exemplo \verb|\textcolor[HTML]{006AD7}|: \textcolor[HTML]{006AD7}{texto azul};
     	\end{itemize}
     	\item Ambiente matem tico inline (comando \verb|$ express o $|): $s = x^2-2x +1$;
     	\item Refer ncia normal (comando \verb|\cite|):
     	\begin{itemize}
     		\item \cite{Agaisse1995};
     		\item \cite{Abedi2014};
     		\item \cite{BtNomenclature2016};
     	\end{itemize}
     	\item Refer ncia normal com mais de uma obra (comando \verb|\cite|):
     	\begin{itemize}
     		\item \cite{Abedi2014, Agaisse1995};
             \item \cite{AgapitoTenfen2014, BtNomenclature2016, Nelson2014};
     	\end{itemize}
     	\item Refer ncia nome e ano (comando \verb|\citeauthorandyear|):
     	\begin{itemize}
     		\item \citeauthorandyear{Agaisse1995};
     		\item \citeauthorandyear{Abedi2014};
     		\item \citeauthorandyear{BtNomenclature2016};
     	\end{itemize}
     \end{itemize}
     
     
     Exemplo 1 de cita o direta:
     
     \begin{citacao}
     	Os 20 amino cidos usualmente encontrados como res duos em prote nas cont m um grupo $\alpha$-carboxil, um grupo $\alpha$-amino e um grupo R distinto substitu do no  tomo de carbono $\alpha$. O  tomo de carbono $\alpha$ de todos os amino cidos, com exce o da glicina,   assim trico e, portanto, os amino cidos podem existir em pelo menos duas formas estereoisom ricas. Somente os estereois meros L, com uma configura o relacionada   configura o absoluta da mol cula de refer ncia L-gliceralde do, s o encontrados em prote nas \cite[p. 81]{Nelson2014}.
     \end{citacao}
     
     Exemplo 2 de cita o direta:
     
     \begin{citacao}
     	\textit{These various insecticidal proteins are synthesized during the stationary phase and accumulate in the mother cell as a crystal inclusion which can account for up to 25\% of the dry weight of the sporulated cells. The amount of crystal protein produced by a B. thuringiensis culture in laboratory conditions (about 0.5 mg of protein per ml) and the size of the crystals (24) indicate that each cell has to synthesize $10^6$ to $2 \times 10^6$ $\delta$-endotoxin molecules during the stationary phase to form a crystal} \cite[p. 1]{Agaisse1995}.
     \end{citacao}
     
     Exemplo de nota de rodap \footnote{Essa   uma nota de rodap !}.
     
     
     \subsubsection{Trabalhos Correlatos}
     
     Pesquise e descreva no m nimo tr s trabalhos correlatos ao seu.
     
     \subsubsubsection{Trabalho 1}
     
     Texto...
     
     \subsubsubsection{Trabalho 2}
     
     Texto...
     
     \subsubsubsection{Trabalho 3}
     
     Texto...
     
     %%%%%%%%%%%%%%%%%%%%%%%%%%%%%%%%%%%%%%%%%%%%%%%%%%%%%%%%%%%%%
     
     \subsection{Metodologia}
     
     Texto da metodologia.
     
     
     
     %%%%%%%%%%%%%%%%%%%%%%%%%%%%%%%%%%%%%%%%%%%%%%%%%%%%%%%%%%%%%
     
     
%%%%%%%%%%%%%%%%%%%%%%%%%%%%%%%%%%%%%%%%%%%%%%%%%%%%%%%%%%%%%%%%%%

	\section{Resultados e Discuss o}
    	
     Texto dos resultados.
     
     
     
%%%%%%%%%%%%%%%%%%%%%%%%%%%%%%%%%%%%%%%%%%%%%%%%%%%%%%%%%%%%%%%%%%

	\section{Conclus es/Conclus es parciais}
	
	Texto das conclus es. Lembre-se de apontar poss veis trabalhos futuros.
	
	\textbf{Obs:} Esta se o deve ser intitulada ``Conclus es Parciais'' em trabalhos de gradua o para a Valida o de Projeto de TCC. Na Avalia o Final de TCC o nome da se o deve ser ``Conclus es''.
    
    
    
%%%%%%%%%%%%%%%%%%%%%%%%%%%%%%%%%%%%%%%%%%%%%%%%%%%%%%%%%%%%%%%%%%

	\section{Cronograma}
	
	Segue abaixo o cronograma das atividades que ser o executadas at  a Avalia o Final de TCC.
	
	\textbf{Obs:} Para facilitar, crie o cronograma usando o modelo do Word contido no projeto (imagens/templateCronograma.docx), ou qualquer outro \textit{software}, salve a imagem e atualize o arquivo imagens/cronograma.png.
	
	\FloatBarrier
	\begin{figure*}[!htbp]
		\centering
		\includegraphics[scale=1]{imagens/cronograma}
	\end{figure*}
	\FloatBarrier
	
	\begin{enumerate}
		\item Descri o da atividade 1;
		\item Descri o da atividade 2;
		\item Descri o da atividade 3;
		\item Descri o da atividade 4;
		\item Descri o da atividade 5.
	\end{enumerate}
	
	\textbf{Obs:} Esta se o deve ser elaborada e estar contida em trabalhos de gradua o para a Valida o de Projeto de TCC. Na Avalia o Final de TCC esta se o n o deve existir, visto que n o haver  atividades ap s a Avalia o Final.
	
	
	% ----------------------------------------------------------
	% ELEMENTOS P S-TEXTUAIS
	% ----------------------------------------------------------
	\postextual
	
	
	%
    % Este   um exemplo da ata de defesa.
    %
    % A ata de defesa ser  fornecida pelo professor orientador ap s a defesa do trabalho
    % e o orientado ser  respons vel em substituir o arquivo de exemplo pelo arquivo final.
    %
    \begin{folhadeaprovacao}
    	\includepdf[pages={1-}]{ataDefesa/exemploAtaDefesa.pdf}
    \end{folhadeaprovacao}
\end{document}